\section{Experimental Setup}

\subsection{Evaluation Tiers}

The benchmark release contains 2{,}500 problems.
The main empirical analysis uses a stratified
500-problem subset with 100 problems from each of the five difficulty
levels.
This larger pool provides substantially tighter uncertainty estimates
for capability curves, account-difficulty analysis, and prompt-sensitivity
measurements.

We also retain a separate paired 100-problem compatibility subset with
20 problems per difficulty level.
That smaller pool extends model coverage to systems currently evaluated
on the 100-problem setting.
The paper therefore distinguishes between a \emph{Large-500} study and a
\emph{Controlled-100} compatibility study, with comparative claims
interpreted within matched pools.

\subsection{Model Coverage}

Completed Large-500 runs are:
DeepSeek-v3.2 (\zs{}/\CoT{}), Gemini~3 Flash (\zs{}/\CoT{}),
Grok~4.1 Fast (\zs{}/\CoT{}), GPT-5.4 Nano (\zs{}/\CoT{}),
Nemotron~3 Super (\zs{}),
Qwen~3.5 Flash (\zs{}/\CoT{}), and Qwen~3.6 Plus (\CoT{}).
These runs are used for the 500-problem figures, account
difficulty table, factor-sensitivity analysis, and paired significance
tests reported in the main text.

Frontier proprietary coverage on the separate Controlled-100 pool
includes GPT-5.2 (\zs{}/\CoT{}) and gpt-oss-120b (\zs{} only).

\subsection{Evaluation and Statistical Protocol}

The evaluation harness decouples model backends from prompting
strategies, enabling controlled comparisons without changing the
benchmark instances or scoring logic.
All runs use deterministic decoding with temperature fixed at \(0\).
A rule-based baseline reaches \fbs{} \(=99.97\).

Where \(p\)-values are cited, they come from paired permutation tests
within a matched evaluation pool.
The main inferential claims use the Large-500 paired tests (\(n=500\));
Controlled-100 results extend the coverage summary for additional
frontier systems.
Cross-pool comparisons are interpreted descriptively.
